\chapter{2004年全国II卷（文）}
\section{单选题}
\begin{problem}
已知集合 \(M = \{x \mid x^2 < 4\}\)，\(N = \{x \mid x^2 - 2x - 3 < 0\}\)，则集合 \(M \cap N =\)（\quad）

\choices
    {\(\{x \mid x < -2\}\)}
    {\(\{x \mid x > 3\}\)}
    {\(\{x \mid -1 < x < 2\}\)}
    {\(\{x \mid 2 < x < 3\}\)}
\end{problem}

\begin{answer}
C
\end{answer}

\begin{solution}
由 \(x^2<4\) 得 \(-2<x<2\)，所以 \(M=(-2,2)\). 又 \(x^2-2x-3=(x-3)(x+1)\)，故 \(-1<x<3\)，即 \(N=(-1,3)\). 因此 \(M\cap N=(-2,2)\cap(-1,3)=(-1,2)\)，所以选 C.
\end{solution}

\begin{problem}
函数 \(y=\dfrac{1}{x+5}\)（\(x\neq-5\)）的反函数是（\quad）

\choices
    {\(y=\dfrac{1}{x}-5\)（\(x\neq0\)）}
    {\(y=x+5\)（\(x\in\R\)）}
    {\(y=\dfrac{1}{x}+5\)（\(x\neq0\)）}
    {\(y=x-5\)（\(x\in\R\)）}
\end{problem}

\begin{answer}
A
\end{answer}

\begin{solution}
由 \(y=\dfrac{1}{x+5}\)，交换 \(x\)，\(y\) 得 \(x=\dfrac{1}{y+5}\). 原函数的值域不包含 \(0\)，所以反函数的定义域为 \(x\neq0\). 解得 \(y=\dfrac{1}{x}-5\)，因此选 A.
\end{solution}

\begin{problem}
曲线 \(y=x^3-3x^2+1\) 在点 \((1,-1)\) 处的切线方程为（\quad）

\choices
    {\(y=3x-4\)}
    {\(y=-3x+2\)}
    {\(y=-4x+3\)}
    {\(y=4x-5\)}
\end{problem}

\begin{answer}
B
\end{answer}

\begin{solution}
设 \(f(x)=x^3-3x^2+1\)，则 \(f'(x)=3x^2-6x\)，从而 \(f'(1)=-3\). 切线经过 \((1,-1)\)，所以其方程为 \(y+1=-3(x-1)\)，即 \(y=-3x+2\)，因此选 B.
\end{solution}

\begin{problem}
已知圆 \(C\) 与圆 \((x-1)^2+y^2=1\) 关于直线 \(y=-x\) 对称，则圆 \(C\) 的方程为（\quad）

\choices
    {\((x+1)^2+y^2=1\)}
    {\(x^2+y^2=1\)}
    {\(x^2+(y+1)^2=1\)}
    {\(x^2+(y-1)^2=1\)}
\end{problem}

\begin{answer}
C
\end{answer}

\begin{solution}
原圆圆心为 \((1,0)\)，半径为 \(1\). 点 \((u,v)\) 关于直线 \(y=-x\) 的对称点为 \((-v,-u)\)，所以圆心 \((1,0)\) 变为 \((0,-1)\)，半径不变. 因此圆 \(C\) 的方程为 \(x^2+(y+1)^2=1\)，所以选 C.
\end{solution}

\begin{problem}
已知函数 \(y=\tan(2x+\varphi)\) 的图象过点 \(\left(\dfrac{\pi}{12},0\right)\)，则 \(\varphi\) 可以是（\quad）

\choices
    {\(-\dfrac{\pi}{6}\)}
    {\(\dfrac{\pi}{6}\)}
    {\(-\dfrac{\pi}{12}\)}
    {\(\dfrac{\pi}{12}\)}
\end{problem}

\begin{answer}
A
\end{answer}

\begin{solution}
将点的横坐标代入，得 \(\tan\left(\dfrac{\pi}{6}+\varphi\right)=0\)，所以 \(\dfrac{\pi}{6}+\varphi=k\pi\)，其中 \(k\in\Z\). 在所给选项中，只有 \(\varphi=-\dfrac{\pi}{6}\) 满足条件，因此选 A.
\end{solution}

\begin{problem}
正四棱锥的侧棱长与底面边长都是 \(1\)，则侧棱与底面所成的角为（\quad）

\choices
    {\(75^{\circ}\)}
    {\(60^{\circ}\)}
    {\(45^{\circ}\)}
    {\(30^{\circ}\)}
\end{problem}

\begin{answer}
C
\end{answer}

\begin{solution}
设正四棱锥顶点为 \(V\)，底面中心为 \(O\)，底面顶点为 \(A\). 由于 \(VO\) 垂直底面，侧棱 \(VA\) 与底面所成的角等于 \(\angle VAO\). 正方形底面边长为 \(1\)，所以 \(AO=\dfrac{\sqrt2}{2}\)，而 \(VA=1\). 在直角三角形 \(VAO\) 中，\(\cos\angle VAO=\dfrac{AO}{VA}=\dfrac{\sqrt2}{2}\)，故 \(\angle VAO=45^{\circ}\)，所以选 C.
\end{solution}

\begin{problem}
函数 \(y=-\e^x\) 的图象（\quad）

\choices
    {与 \(y=\e^x\) 的图象关于 \(y\) 轴对称}
    {与 \(y=\e^x\) 的图象关于坐标原点对称}
    {与 \(y=\e^{-x}\) 的图象关于 \(y\) 轴对称}
    {与 \(y=\e^{-x}\) 的图象关于坐标原点对称}
\end{problem}

\begin{answer}
D
\end{answer}

\begin{solution}
曲线 \(y=\e^{-x}\) 上的点 \((x,\e^{-x})\) 关于坐标原点对称后变为 \((-x,-\e^{-x})\). 令新横坐标为 \(X=-x\)，则新纵坐标为 \(-\e^X\)，所以对称后的曲线为 \(y=-\e^x\)，因此选 D.
\end{solution}

\begin{problem}
已知点 \(A(1,2)\)，\(B(3,1)\)，则线段 \(AB\) 的垂直平分线的方程为（\quad）

\choices
    {\(4x+2y=5\)}
    {\(4x-2y=5\)}
    {\(x+2y=5\)}
    {\(x-2y=5\)}
\end{problem}

\begin{answer}
B
\end{answer}

\begin{solution}
线段 \(AB\) 的中点为 \(\left(2,\dfrac{3}{2}\right)\)，且 \(AB\) 的斜率为 \(\dfrac{1-2}{3-1}=-\dfrac{1}{2}\). 因此垂直平分线的斜率为 \(2\)，其方程为 \(y-\dfrac{3}{2}=2(x-2)\)，化简得 \(4x-2y=5\)，所以选 B.
\end{solution}

\begin{problem}
已知向量 \(\bs{a}\)，\(\bs{b}\) 满足：\(\left|\bs{a}\right|=1\)，\(\left|\bs{b}\right|=2\)，\(\left|\bs{a}-\bs{b}\right|=2\)，则 \(\left|\bs{a}+\bs{b}\right|=\)（\quad）

\choices
    {\(1\)}
    {\(\sqrt2\)}
    {\(\sqrt5\)}
    {\(\sqrt6\)}
\end{problem}

\begin{answer}
D
\end{answer}

\begin{solution}
由 \(\left|\bs{a}-\bs{b}\right|^2=\left|\bs{a}\right|^2+\left|\bs{b}\right|^2-2\bs{a}\cdot\bs{b}\)，得 \(4=1+4-2\bs{a}\cdot\bs{b}\)，所以 \(\bs{a}\cdot\bs{b}=\dfrac{1}{2}\). 因而 \(\left|\bs{a}+\bs{b}\right|^2=\left|\bs{a}\right|^2+\left|\bs{b}\right|^2+2\bs{a}\cdot\bs{b}=1+4+1=6\)，从而 \(\left|\bs{a}+\bs{b}\right|=\sqrt6\)，所以选 D.
\end{solution}

\begin{problem}
已知球 \(O\) 的半径为 \(1\)，\(A\)，\(B\)，\(C\) 三点都在球面上，且每两点间的球面距离为 \(\dfrac{\pi}{2}\)，则球心 \(O\) 到平面 \(ABC\) 的距离为（\quad）

\choices
    {\(\dfrac{1}{3}\)}
    {\(\dfrac{\sqrt3}{3}\)}
    {\(\dfrac{2}{3}\)}
    {\(\dfrac{\sqrt6}{3}\)}
\end{problem}

\begin{answer}
B
\end{answer}

\begin{solution}
球面距离等于相应圆心角乘半径，故 \(\angle AOB=\angle BOC=\angle COA=\dfrac{\pi}{2}\). 取 \(\bs{u}=\overrightarrow{OA}\)，\(\bs{v}=\overrightarrow{OB}\)，\(\bs{w}=\overrightarrow{OC}\)，则这三个向量两两垂直且模长均为 \(1\). 令 \(\bs{n}=\bs{u}+\bs{v}+\bs{w}\)，则 \(\bs{n}\cdot\bs{u}=\bs{n}\cdot\bs{v}=\bs{n}\cdot\bs{w}=1\)，所以 \(\bs{n}\) 是平面 \(ABC\) 的法向量，且 \(\left|\bs{n}\right|=\sqrt3\). 平面 \(ABC\) 的方程可写为 \(\bs{n}\cdot\overrightarrow{OX}=1\)，因此球心到该平面的距离为 \(\dfrac{1}{\left|\bs{n}\right|}=\dfrac{\sqrt3}{3}\)，所以选 B.
\end{solution}

\begin{problem}
函数 \(y=\sin^4x+\cos^2x\) 的最小正周期为（\quad）

\choices
    {\(\dfrac{\pi}{4}\)}
    {\(\dfrac{\pi}{2}\)}
    {\(\pi\)}
    {\(2\pi\)}
\end{problem}

\begin{answer}
B
\end{answer}

\begin{solution}
利用 \(\sin^2x+\cos^2x=1\)，有 \(\sin^4x+\cos^2x=1-\sin^2x\cos^2x=\dfrac{7}{8}+\dfrac{1}{8}\cos4x\). 因此 \(\dfrac{\pi}{2}\) 是一个正周期. 若 \(T>0\) 是周期，则 \(\cos4(x+T)=\cos4x\) 对任意 \(x\) 成立，从而 \(4T=2k\pi\)，其中 \(k\in\Z\). 所以最小正周期为 \(\dfrac{\pi}{2}\)，故选 B.
\end{solution}

\begin{problem}
在由数字 \(1\)，\(2\)，\(3\)，\(4\)，\(5\) 组成的所有没有重复数字的 \(5\) 位数中，大于 \(23145\) 且小于 \(43521\) 的数共有（\quad）

\choices
    {\(56\text{ 个}\)}
    {\(57\text{ 个}\)}
    {\(58\text{ 个}\)}
    {\(60\text{ 个}\)}
\end{problem}

\begin{answer}
C
\end{answer}

\begin{solution}
按首位分类计数. 首位为 \(2\) 且大于 \(23145\) 时，第二位为 \(4\) 或 \(5\) 的有 \(2\cdot3!=12\) 个；第二位为 \(3\) 且第三位为 \(4\) 或 \(5\) 的有 \(2\cdot2!=4\) 个；前缀为 \(231\) 且第四位为 \(5\) 的有 \(1\) 个，共 \(17\) 个.

首位为 \(3\) 的数全部满足不等式，共有 \(4!=24\) 个. 首位为 \(4\) 且小于 \(43521\) 时，第二位为 \(1\) 或 \(2\) 的有 \(2\cdot3!=12\) 个；前两位为 \(43\) 且第三位为 \(1\) 或 \(2\) 的有 \(2\cdot2!=4\) 个；前三位为 \(435\) 且第四位为 \(1\) 的有 \(1\) 个，共 \(17\) 个. 因此总数为 \(17+24+17=58\)，所以选 C.
\end{solution}

\section{填空题}
\begin{problem}
已知 \(a\) 为实数，\((x+a)^{10}\) 展开式中 \(x^7\) 的系数是 \(-15\)，则 \(a=\)\fillinblank{}.
\end{problem}

\begin{answer}
\(-\dfrac{1}{2}\)
\end{answer}

\begin{solution}
展开式中 \(x^7\) 项由 \(7\) 个 \(x\) 和 \(3\) 个 \(a\) 相乘得到，所以其系数为 \(\mathrm C_{10}^{3}a^3=120a^3\). 由题意，\(120a^3=-15\)，得 \(a^3=-\dfrac{1}{8}\). 因为 \(a\) 为实数，所以 \(a=-\dfrac{1}{2}\).
\end{solution}

\begin{problem}
设 \(x\)，\(y\) 满足约束条件
\(\begin{cases}
x\geqslant0,\\
x\geqslant y,\\
2x-y\leqslant1,
\end{cases}\)
则 \(z=3x+2y\) 的最大值是\fillinblank{}.
\end{problem}

\begin{answer}
5
\end{answer}

\begin{solution}
由约束条件得 \(y\leqslant x\) 且 \(y\geqslant2x-1\). 要使这样的 \(y\) 存在，必须有 \(2x-1\leqslant x\)，即 \(x\leqslant1\). 结合 \(x\geqslant0\)，得 \(0\leqslant x\leqslant1\). 对固定的 \(x\)，\(z=3x+2y\) 随 \(y\) 增大而增大，因此取 \(y=x\) 时最大，此时 \(z=5x\leqslant5\). 当 \((x,y)=(1,1)\) 时取等号，所以最大值为 \(5\).
\end{solution}

\begin{problem}
设中心在原点的椭圆与双曲线 \(2x^2-2y^2=1\) 有公共的焦点，且它们的离心率互为倒数，则该椭圆的方程是\fillinblank{}.
\end{problem}

\begin{answer}
\(\dfrac{x^2}{2}+y^2=1\)
\end{answer}

\begin{solution}
双曲线方程化为 \(\dfrac{x^2}{1/2}-\dfrac{y^2}{1/2}=1\)，所以 \(a_{\mathrm h}^2=b_{\mathrm h}^2=\dfrac{1}{2}\)，且 \(c_{\mathrm h}^2=a_{\mathrm h}^2+b_{\mathrm h}^2=1\). 因而 \(c_{\mathrm h}=1\)，双曲线的离心率为 \(e_{\mathrm h}=\dfrac{c_{\mathrm h}}{a_{\mathrm h}}=\sqrt2\).

设椭圆的半长轴、半短轴和焦半距分别为 \(a_{\mathrm e}\)，\(b_{\mathrm e}\)，\(c_{\mathrm e}\). 由公共焦点知 \(c_{\mathrm e}=1\)，由离心率互为倒数知 \(e_{\mathrm e}=\dfrac{1}{\sqrt2}=\dfrac{c_{\mathrm e}}{a_{\mathrm e}}\)，所以 \(a_{\mathrm e}=\sqrt2\). 于是 \(b_{\mathrm e}^2=a_{\mathrm e}^2-c_{\mathrm e}^2=2-1=1\)，故椭圆方程为 \(\dfrac{x^2}{2}+y^2=1\).
\end{solution}

\begin{problem}
下面是关于四棱柱的四个命题：

\begin{circlelist}
\item 若有两个侧面垂直于底面，则该四棱柱为直四棱柱；
\item 若两个过相对侧棱的截面都垂直于底面，则该四棱柱为直四棱柱；
\item 若四个侧面两两全等，则该四棱柱为直四棱柱；
\item 若四棱柱的四条对角线两两相等，则该四棱柱为直四棱柱.
\end{circlelist}

其中，真命题的编号是\fillinblank{}.（写出所有真命题的编号）
\end{problem}

\begin{answer}
2、4
\end{answer}

\begin{solution}
设四棱柱的底面在水平面内，侧棱向量为 \(\bs{v}\)，其在底面内的投影为 \(\bs{p}\).

对于命题 1，取底面为正方形，侧棱向量取 \((1,0,1)\). 过底面中两条相对且平行的边的两个侧面都是竖直平面，因而都垂直于底面，但侧棱不垂直于底面，所以该四棱柱不是直四棱柱，命题 1 是假命题.

对于命题 2，过相对侧棱的两个截面分别包含底面对角线 \(AC\)，\(BD\) 和侧棱方向. 截面垂直于底面时，\(\bs{p}\) 必分别平行于 \(AC\) 和 \(BD\). 由于非退化四边形的两条对角线不平行，只能有 \(\bs{p}=\bs{0}\)，即侧棱垂直于底面，所以命题 2 是真命题.

对于命题 3，取底面为单位正方形，侧棱向量取 \((1,1,1)\). 每个侧面都是由一个单位边向量和向量 \((1,1,1)\) 张成的平行四边形；这两个向量的长度分别为 \(1\)，\(\sqrt3\)，且相应内积的绝对值都为 \(1\)，所以四个侧面两两全等. 但侧棱不垂直于底面，命题 3 是假命题.

对于命题 4，设底面两条对角线的向量分别为 \(\bs{u}=\overrightarrow{AC}\)，\(\bs{w}=\overrightarrow{BD}\). 四条对角线的向量可分别写成 \(\bs{v}+\bs{u}\)，\(\bs{v}-\bs{u}\)，\(\bs{v}+\bs{w}\)，\(\bs{v}-\bs{w}\). 由四条对角线两两相等，特别有 \(\left|\bs{v}+\bs{u}\right|=\left|\bs{v}-\bs{u}\right|\) 和 \(\left|\bs{v}+\bs{w}\right|=\left|\bs{v}-\bs{w}\right|\). 两式平方后分别得到 \(\bs{v}\cdot\bs{u}=0\) 和 \(\bs{v}\cdot\bs{w}=0\). 两条对角线张成底面，所以 \(\bs{v}\) 垂直于底面，命题 4 是真命题. 因此真命题的编号为 2、4.
\end{solution}

\section{解答题}
\begin{problem}
已知等差数列 \(\{a_n\}\)，\(a_2=9\)，\(a_5=21\).

\begin{enumerate}
\item 求 \(\{a_n\}\) 的通项公式；
\item 令 \(b_n=2^{a_n}\)，求数列 \(\{b_n\}\) 的前 \(n\) 项和 \(S_n\).
\end{enumerate}
\end{problem}

\begin{solution}
设等差数列的首项为 \(a_1\)，公差为 \(d\). 由已知得
\[
\begin{cases}
a_1+d=9,\\
a_1+4d=21.
\end{cases}
\]
两式相减得 \(3d=12\)，所以 \(d=4\)，再由 \(a_1+d=9\) 得 \(a_1=5\). 因此 \(a_n=a_1+(n-1)d=5+4(n-1)=4n+1\).

于是 \(b_n=2^{a_n}=2^{4n+1}=32\cdot16^{n-1}\)，这是首项为 \(32\)、公比为 \(16\) 的等比数列. 所以 \(S_n=32\dfrac{16^n-1}{16-1}=\dfrac{32(16^n-1)}{15}\).
\end{solution}

\begin{problem}
已知锐角三角形 \(ABC\) 中，\(\sin(A+B)=\dfrac{3}{5}\)，\(\sin(A-B)=\dfrac{1}{5}\).

\begin{enumerate}
\item 求证：\(\tan A=2\tan B\)；
\item 设 \(AB=3\)，求 \(AB\) 边上的高.
\end{enumerate}
\end{problem}

\begin{solution}
因为三角形 \(ABC\) 为锐角三角形，所以 \(A\)，\(B\)，\(C\) 均为锐角，且 \(A+B=\pi-C>\dfrac{\pi}{2}\). 又 \(\sin(A+B)=\sin C=\dfrac{3}{5}\)，故 \(\cos C=\dfrac{4}{5}\)，从而 \(\cos(A+B)=-\cos C=-\dfrac{4}{5}\).

记 \(X=\sin A\cos B\)，\(Y=\cos A\sin B\). 由和差公式得 \(X+Y=\dfrac{3}{5}\)，\(X-Y=\dfrac{1}{5}\)，所以 \(X=\dfrac{2}{5}\)，\(Y=\dfrac{1}{5}\). 因为 \(A\)，\(B\) 为锐角，\(\cos A\cos B>0\)，于是 \(\dfrac{\tan A}{\tan B}=\dfrac{\sin A\cos B}{\cos A\sin B}=\dfrac{X}{Y}=2\)，即 \(\tan A=2\tan B\).

由 \(\sin(A-B)=\dfrac{1}{5}>0\) 且 \(A-B\in\left(-\dfrac{\pi}{2},\dfrac{\pi}{2}\right)\)，得 \(A>B\)，从而 \(\cos(A-B)=\sqrt{1-\dfrac{1}{25}}=\dfrac{2\sqrt6}{5}\). 因此 \(\sin A\sin B=\dfrac{\cos(A-B)-\cos(A+B)}{2}=\dfrac{\frac{2\sqrt6}{5}+\frac{4}{5}}{2}=\dfrac{\sqrt6+2}{5}\).

设三角形外接圆半径为 \(R\). 由正弦定理，\(AB=2R\sin C=3\)，故 \(2R=5\). 设 \(AB\) 边上的高为 \(h\)，由三角形面积的两种表示方法，\(\dfrac{1}{2}AB\cdot h=\dfrac{1}{2}AC\cdot BC\sin C\). 再由正弦定理 \(AC=2R\sin B\)，\(BC=2R\sin A\)，得 \(h=2R\sin A\sin B=5\cdot\dfrac{\sqrt6+2}{5}=2+\sqrt6\).
\end{solution}

\begin{problem}
已知 \(8\text{ 支}\) 球队中有 \(3\text{ 支}\) 弱队，以抽签方式将这 \(8\text{ 支}\) 球队分为 \(A\)，\(B\) 两组，每组 \(4\text{ 支}\)，求：

\begin{enumerate}
\item \(A\)，\(B\) 两组中有一组恰有 \(2\text{ 支}\) 弱队的概率；
\item \(A\) 组中至少有两支弱队的概率.
\end{enumerate}
\end{problem}

\begin{solution}
从 \(8\) 支球队中选出进入 \(A\) 组的 \(4\) 支，基本事件总数为 \(\mathrm C_8^4=70\). 设 \(A\) 组中有 \(k\) 支弱队，则 \(B\) 组中有 \(3-k\) 支弱队.

第（1）问中，恰有一组有 \(2\) 支弱队等价于 \(k=1\) 或 \(k=2\)，所以所求概率为 \(\dfrac{\mathrm C_3^1\mathrm C_5^3+\mathrm C_3^2\mathrm C_5^2}{\mathrm C_8^4}=\dfrac{3\cdot10+3\cdot10}{70}=\dfrac{6}{7}\).

第（2）问中，\(A\) 组至少有两支弱队等价于 \(k=2\) 或 \(k=3\)，所以所求概率为 \(\dfrac{\mathrm C_3^2\mathrm C_5^2+\mathrm C_3^3\mathrm C_5^1}{\mathrm C_8^4}=\dfrac{3\cdot10+1\cdot5}{70}=\dfrac{1}{2}\).
\end{solution}

\begin{problem}
如图，直三棱柱 \(ABC-A_1B_1C_1\) 中，\(\angle ACB=90^{\circ}\)，\(AC=1\)，\(CB=\sqrt2\)，侧棱 \(AA_1=1\)，侧面 \(AA_1B_1B\) 的两条对角线交点为 \(D\)，\(B_1C_1\) 的中点为 \(M\).

\begin{enumerate}
\item 求证：\(CD\perp\) 平面 \(BDM\)；
\item 求面 \(B_1BD\) 与面 \(CBD\) 所成二面角的大小.
\end{enumerate}

\bitmapfigure[max width=0.42\linewidth]{img/2004/national_paper_2_liberal/q18_fig1.png}
\end{problem}

\begin{solution}
建立空间直角坐标系，使 \(C\) 为原点，\(CA\)，\(CB\)，\(CC_1\) 分别为三个坐标轴的正方向. 则 \(C=(0,0,0)\)，\(A=(1,0,0)\)，\(B=(0,\sqrt2,0)\)，\(A_1=(1,0,1)\)，\(B_1=(0,\sqrt2,1)\)，\(C_1=(0,0,1)\).

因为 \(D\) 是矩形 \(AA_1B_1B\) 的两条对角线交点，且 \(M\) 是 \(B_1C_1\) 的中点，所以 \(D=\left(\dfrac{1}{2},\dfrac{\sqrt2}{2},\dfrac{1}{2}\right)\)，\(M=\left(0,\dfrac{\sqrt2}{2},1\right)\). 于是 \(\overrightarrow{CD}=\left(\dfrac{1}{2},\dfrac{\sqrt2}{2},\dfrac{1}{2}\right)\)，\(\overrightarrow{DB}=\left(-\dfrac{1}{2},\dfrac{\sqrt2}{2},-\dfrac{1}{2}\right)\)，\(\overrightarrow{DM}=\left(-\dfrac{1}{2},0,\dfrac{1}{2}\right)\).

直接计算得 \(\overrightarrow{CD}\cdot\overrightarrow{DB}=0\)，\(\overrightarrow{CD}\cdot\overrightarrow{DM}=0\). 直线 \(DB\)，\(DM\) 是平面 \(BDM\) 内相交的两条直线，所以 \(CD\perp\) 平面 \(BDM\).

求二面角时，公共棱为 \(BD\). 取 \(\bs{e}=\overrightarrow{DB}=\left(-\dfrac{1}{2},\dfrac{\sqrt2}{2},-\dfrac{1}{2}\right)\)，则 \(\left|\bs{e}\right|=1\). 在平面 \(B_1BD\) 内，\(\overrightarrow{DB_1}=\left(-\dfrac{1}{2},\dfrac{\sqrt2}{2},\dfrac{1}{2}\right)\)，且 \(\overrightarrow{DB_1}\cdot\bs{e}=\dfrac{1}{2}\). 将其减去在 \(BD\) 上的投影，得到垂直于 \(BD\) 的向量 \(\bs{u}=\overrightarrow{DB_1}-\dfrac{1}{2}\bs{e}=\left(-\dfrac{1}{4},\dfrac{\sqrt2}{4},\dfrac{3}{4}\right)\)，且 \(\left|\bs{u}\right|=\dfrac{\sqrt3}{2}\).

在平面 \(CBD\) 内，\(\overrightarrow{DC}=\left(-\dfrac{1}{2},-\dfrac{\sqrt2}{2},-\dfrac{1}{2}\right)\) 与 \(\bs{e}\) 垂直，且 \(\left|\overrightarrow{DC}\right|=1\). 因此所求二面角的平面角为 \(\theta\)，满足 \(\cos\theta=\dfrac{\bs{u}\cdot\overrightarrow{DC}}{\left|\bs{u}\right|\left|\overrightarrow{DC}\right|}=-\dfrac{1/2}{(\sqrt3/2)\cdot1}=-\dfrac{\sqrt3}{3}\). 所以 \(\theta=\arccos\left(-\dfrac{\sqrt3}{3}\right)=\pi-\arccos\dfrac{\sqrt3}{3}\).
\end{solution}

\begin{problem}
若函数 \(f(x)=\dfrac{1}{3}x^3-\dfrac{1}{2}ax^2+(a-1)x+1\) 在区间 \((1,4)\) 内为减函数，在区间 \((6,+\infty)\) 上为增函数，试求实数 \(a\) 的取值范围.
\end{problem}

\begin{solution}
求导得 \(f'(x)=x^2-ax+a-1=(x-1)(x-a+1)\).

当 \(x\in(1,4)\) 时，\(x-1>0\). 要使函数在该区间内为减函数，需使 \(x-a+1<0\) 对任意 \(x\in(1,4)\) 成立，这等价于 \(a\geqslant5\). 当 \(a=5\) 时，\(x-a+1=x-4<0\) 对任意 \(1<x<4\) 成立.

当 \(x\in(6,+\infty)\) 时，仍有 \(x-1>0\). 要使函数在该区间内为增函数，需使 \(x-a+1>0\) 对任意 \(x>6\) 成立，这等价于 \(a\leqslant7\). 当 \(a=7\) 时，\(x-a+1=x-6>0\) 对任意 \(x>6\) 成立.

综上，实数 \(a\) 的取值范围为 \([5,7]\).
\end{solution}

\begin{problem}
给定抛物线 \(C:y^2=4x\)，\(F\) 是 \(C\) 的焦点，过点 \(F\) 的直线 \(l\) 与 \(C\) 相交于 \(A\)，\(B\) 两点.

\begin{enumerate}
\item 设 \(l\) 的斜率为 \(1\)，求 \(\overrightarrow{OA}\) 与 \(\overrightarrow{OB}\) 夹角的大小；
\item 设 \(\overrightarrow{FB}=\lambda\overrightarrow{AF}\)，若 \(\lambda\in[4,9]\)，求 \(l\) 在 \(y\) 轴上截距的变化范围.
\end{enumerate}
\end{problem}

\begin{solution}
抛物线 \(y^2=4x\) 的焦点为 \(F=(1,0)\).

当直线斜率为 \(1\) 时，直线方程为 \(y=x-1\). 代入抛物线方程得 \((x-1)^2=4x\)，即 \(x^2-6x+1=0\). 所以两个交点可取为 \(A=(3-2\sqrt2,2-2\sqrt2)\)，\(B=(3+2\sqrt2,2+2\sqrt2)\).

于是 \(\overrightarrow{OA}\cdot\overrightarrow{OB}=(3-2\sqrt2)(3+2\sqrt2)+(2-2\sqrt2)(2+2\sqrt2)=1-4=-3\). 又 \(\left|\overrightarrow{OA}\right|^2=29-20\sqrt2\)，\(\left|\overrightarrow{OB}\right|^2=29+20\sqrt2\)，故 \(\left|\overrightarrow{OA}\right|^2\left|\overrightarrow{OB}\right|^2=41\). 设所求夹角为 \(\theta\)，则 \(\cos\theta=-\dfrac{3}{\sqrt{41}}\)，所以 \(\theta=\arccos\left(-\dfrac{3}{\sqrt{41}}\right)\).

设 \(\overrightarrow{AF}=\bs{u}=(p,q)\). 由题设 \(\overrightarrow{FB}=\lambda\bs{u}\)，且 \(F=(1,0)\)，故 \(A=(1-p,-q)\)，\(B=(1+\lambda p,\lambda q)\). 将两点分别代入 \(y^2=4x\)，得 \(q^2=4(1-p)\) 和 \(\lambda^2q^2=4(1+\lambda p)\). 消去 \(q^2\)，得到 \(\lambda^2(1-p)=1+\lambda p\)，从而 \(p=\dfrac{\lambda-1}{\lambda}\)，\(q^2=\dfrac{4}{\lambda}\).

由于 \(\lambda\in[4,9]\)，有 \(p>0\)，直线 \(l\) 的斜率为 \(k=\dfrac{q}{p}=\pm\dfrac{2\sqrt{\lambda}}{\lambda-1}\). 直线过 \(F=(1,0)\)，所以方程为 \(y=k(x-1)\)，在 \(y\) 轴上的截距为 \(b=-k=\pm\dfrac{2\sqrt{\lambda}}{\lambda-1}\).

令 \(t=\sqrt{\lambda}\)，则 \(2\leqslant t\leqslant3\)，且 \(g(t)=\dfrac{2t}{t^2-1}\)，\(g'(t)=-\dfrac{2(t^2+1)}{(t^2-1)^2}<0\). 因此 \(\dfrac{3}{4}=g(3)\leqslant g(t)\leqslant g(2)=\dfrac{4}{3}\)，所以截距的变化范围为 \(\left[-\dfrac{4}{3},-\dfrac{3}{4}\right]\cup\left[\dfrac{3}{4},\dfrac{4}{3}\right]\).
\end{solution}
